
\begin{theorem}[Parent-Hamiltonian domination with \texorpdfstring{$\eta=0$}{eta=0}]
\label{lem:norm-stable-gap-parent-path}
\label{thm:eta-zero-parent-domination}
Fix finite volume $V$ and let $T_V$ be the primitive blocked admissible transfer operator with
Perron vector $v_V^\ast$. Let $H_{\mathrm{par}}(V)$ be the blocked admissible parent Hamiltonian of
\Cref{lem:parent-gap-from-transfer-block}, with
\[
m_0(V):=\operatorname{gap}(H_{\mathrm{par}}(V))>0.
\]
Let the physical admissible Hamiltonian be written as
\[
H_{\mathrm{adm}}(V)
=
H_{\mathrm{par}}(V)+K_\Sigma(V)+K_c(V)+K_\Theta(V)+K_g(V)+H_{\mathrm{tr,pos}}(V),
\]
where each added term is positive semidefinite and vanishes on the Perron vacuum line generated by
$v_V^\ast$. Then
\[
H_{\mathrm{adm}}(V)\ge H_{\mathrm{par}}(V).
\]
Consequently the admissible Hamiltonian has the same vacuum line and satisfies
\[
\operatorname{gap}(H_{\mathrm{adm}}(V))
\ge
\operatorname{gap}(H_{\mathrm{par}}(V))
=
m_0(V).
\]
\end{theorem}

\begin{proof}
Every summand in
\[
H_{\mathrm{adm}}(V)-H_{\mathrm{par}}(V)
\]
is positive semidefinite by construction and annihilates the Perron vacuum line, hence the operator
inequality is immediate and the ground state is unchanged. The gap bound then follows from the
min-max principle on the orthogonal complement of that line.
\end{proof}

\begin{theorem}[Block 4: Stability of the admissible mass gap under the thermodynamic limit]
\label{thm:gap-stability-thermo-limit}
Let $(V_n)_{n\ge 1}$ be an admissible exhaustion, let $H_{\mathrm{par}}(V_n)$ be the blocked
parent Hamiltonians of \Cref{lem:parent-gap-from-transfer-block}, and let the uniform parent-gap
bound of \Cref{thm:uniform-parent-gap-injective} be
\[
\inf_n \operatorname{gap}(H_{\mathrm{par}}(V_n))=:m_0>0.
\]
Assume moreover that each finite-volume admissible Hamiltonian $H_{\mathrm{adm}}(V_n)$ satisfies
the parent-domination hypothesis of \Cref{thm:eta-zero-parent-domination}. Then the unique vacuum
persists, the spectral gap of the admissible Hamiltonian remains strictly positive in the
thermodynamic limit, and admissible
correlators cluster exponentially.
\end{theorem}

\begin{proof}
By \Cref{thm:eta-zero-parent-domination}, every finite-volume admissible Hamiltonian obeys
\[
\operatorname{gap}(H_{\mathrm{adm}}(V_n))
\ge
\operatorname{gap}(H_{\mathrm{par}}(V_n))
\ge
m_0>0.
\]
Passing to the thermodynamic limit of \Cref{thm:thermodynamic-limit} preserves the unique vacuum,
and standard gapped quasilocal dynamics then yields exponential clustering.
\end{proof}

\begin{remark}[Closure logic of Blocks 2--4]
\label{rem:vacuum-honest-status}
\TFPTxref{prop:strict-positive-c6-green,cor:uniform-positive-c6-green} provide the local transport
positivity input. Block 2 upgrades that local positivity to a primitive transfer theorem in finite
volume, Block 3 passes the Perron vacuum to the thermodynamic limit, the blocked-transfer theorem
upgrades strict positivity plus diagonal separation to irreducibility, full matrix algebra, and
therefore injectivity, and Block 4 carries the mass gap through explicit parent domination with
$\eta=0$. This closes the vacuum sector inside the theorem class and removes the former
residual-status tag attached to $F_{\mathrm{vac}}$.
\end{remark}

\begin{theorem}[Nonperturbative admissible QFT closure]
\label{thm:nonperturbative-admissible-qft}
Let
\[
m_0:=\inf_n \operatorname{gap}(H_{\mathrm{par}}(V_n))>0
\]
be the uniform parent gap provided by \Cref{thm:uniform-parent-gap-injective}. Define
\[
m_*:=\min\left\{m_0,\,m_-,\,m_3\right\}.
\]
Under \Cref{thm:cp-even-vacuum-uniqueness,thm:thermodynamic-limit,thm:gap-stability-thermo-limit}
together with the transport gap estimate of \TFPTxref{prop:gap-lift-full-operator} and the positivity
bounds of \TFPTxref{thm:callias-higgs-selection}, the admissible branch defines a nonperturbative
gapped QFT closure. After removing the exact gauge zero modes, the reconstructed admissible
Hamiltonian satisfies
\[
\operatorname{spec}(H_{\mathrm{adm}})\cap(0,m_*)=\varnothing.
\]
Hence connected admissible Euclidean correlators obey exponential clustering:
\[
\bigl|\langle O_1(x)O_2(0)\rangle_c\bigr|
\le
C_{12}e^{-m_*|x|}.
\]
\end{theorem}

\begin{proof}[Proof sketch]
The primitive transfer theorem fixes the unique CP-even vacuum. The thermodynamic limit then
produces a unique quasilocal limit state. The gap-stability theorem preserves a strictly positive
spectral gap on that branch through explicit operator domination over the parent Hamiltonian with
$\eta=0$ rather than through a small-norm perturbative corridor. The transport and bosonic
positivity bounds identify the common lower scale $m_*$ for the admissible sector. Reflection
positivity and Lorentzian reconstruction then convert the Euclidean closure data into a Lorentzian
gapped Hamiltonian, and exponential clustering follows on the admissible branch. The decisive
finite-block input is now the chain
\[
\text{strict positivity}\Rightarrow\text{irreducibility}\Rightarrow\text{full matrix algebra}
\Rightarrow\text{injectivity}\Rightarrow\text{uniform gap},
\]
not a separate injectivity assumption.
\end{proof}

\subsection{Admissible Schwinger functions, Minkowski locality, and scattering closure}

\begin{definition}[Laboratory scattering sector]
\label{def:laboratory-scattering-sector}
Let $\mathcal U_{\mathrm{scat}}\subset M$ denote a low-curvature branch region on which
\[
\|g_{\mu\nu}-\eta_{\mu\nu}\|_{C^2(\mathcal U_{\mathrm{scat}})}\ll 1,
\qquad
\|\nabla\chigeo\|_{C^0(\mathcal U_{\mathrm{scat}})}\ll m_*\chigeozero,
\]
and
\[
\left\|
\Phi-\frac{1}{\sqrt2}\binom{0}{\vgeo}
\right\|_{C^0(\mathcal U_{\mathrm{scat}})}
\ll
\vgeo.
\]
All scattering statements below are formulated on $\mathcal U_{\mathrm{scat}}$.
\end{definition}

\begin{definition}[Admissible Schwinger functions]
\label{def:admissible-schwinger-functions}
Let
\[
W_{\mathrm{rel}}[J,\eta,\bar\eta]:=\log Z_{\mathrm{rel}}[J,\eta,\bar\eta].
\]
For local source insertions supported in the Euclidean continuation of
$\mathcal U_{\mathrm{scat}}$, define the truncated admissible Schwinger distributions by
\[
S_n^T(x_1,\ldots,x_n)
:=
\left.
\frac{\delta^n W_{\mathrm{rel}}}{\delta\mathcal J(x_1)\cdots\delta\mathcal J(x_n)}
\right|_{\mathcal J=0},
\qquad
\mathcal J:=(J,\eta,\bar\eta).
\]
\end{definition}

\begin{theorem}[\textnormal{[M]} Full admissible Osterwalder--Schrader closure]
\label{thm:full-admissible-os-closure}
On the Euclidean continuation of $\mathcal U_{\mathrm{scat}}$, the family
$\{S_n^T\}_{n\ge 1}$ satisfies the local Osterwalder--Schrader axioms: regularity, Euclidean
covariance, reflection positivity, symmetry, and clustering. Hence there exists a reconstructed
Wightman theory
\[
(\mathcal H_{\mathrm{OS}},U(a,\Lambda),\Omega,\{\Phi_i\})
\]
whose Schwinger functions are the boundary values of the admissible Euclidean family
$\{S_n^T\}_{n\ge 1}$.
\end{theorem}

\begin{proof}
Regularity and moment bounds are recorded in
\Cref{lem:tempered-schwinger-bounds}, while graded Euclidean covariance and symmetry on $\Padm$ are
recorded in \Cref{lem:graded-euclidean-covariance-padm}. Reflection positivity is exactly
\TFPTxref{thm:reflection-positivity-admissible}, and clustering is the exponential clustering statement
of \TFPTxref{thm:nonperturbative-admissible-qft}. Therefore the admissible Schwinger family satisfies
the hypotheses of the tailored OS/CAR reconstruction theorem
\Cref{thm:tailored-os-car-reconstruction}, which yields the claimed reconstructed Wightman theory.
\end{proof}

\begin{theorem}[Admissible local Minkowski net and microcausality]
\label{thm:admissible-local-minkowski-net}
Let $\mathfrak A_{\mathrm{adm}}(\mathcal O)$ be the physical admissible local net of
\TFPTxref{def:physical-admissible-local-net}. Then:
\begin{enumerate}[label=\textup{(\roman*)}, leftmargin=1.8em, topsep=3pt, itemsep=2pt]
    \item if $\mathcal O_1\subset\mathcal O_2$, then
    \[
    \mathfrak A_{\mathrm{adm}}(\mathcal O_1)\subset \mathfrak A_{\mathrm{adm}}(\mathcal O_2);
    \]
    \item if $\mathcal O_1\perp\mathcal O_2$, then
    \[
    [A_1,A_2]=0
    \qquad
    (A_1\in\mathfrak A_{\mathrm{adm}}(\mathcal O_1),\ A_2\in\mathfrak A_{\mathrm{adm}}(\mathcal O_2)).
    \]
\end{enumerate}
Equivalently, the admissible branch reconstructs a causal Minkowski local net on
$\mathcal H_{\mathrm{adm}}$.
\end{theorem}

\begin{proof}
Isotony is immediate from the support definition of the unreduced local net
$\mathfrak A_M(\mathcal O)$. For locality, consider the Wick-rotated quadratic operators on the
low-curvature branch. Scalar and bosonic modes are normally hyperbolic, fermions are
Dirac-hyperbolic, and gauge fields become hyperbolic after covariant gauge fixing and BRST
reduction. Hence the advanced and retarded fundamental solutions $E^\pm$ exist and satisfy
\[
\operatorname{supp}(E^\pm)\subset J^\pm.
\]
Therefore the Pauli--Jordan propagator
\[
E:=E^+-E^-
\]
has support in the causal cone, so for spacelike separated test-function supports one has
\[
[\Phi_i(f),\Phi_j(g)]_{\mp}=0.
\]
Now let $A_1$ and $A_2$ be local operators preserving $\mathcal H_{\mathrm{adm}}$. Since they
commute on $\mathcal H_{\mathrm{OS}}$ for spacelike separated supports, their restrictions to the
invariant subspace $\mathcal H_{\mathrm{adm}}$ also commute:
\[
[A_1|_{\mathcal H_{\mathrm{adm}}},A_2|_{\mathcal H_{\mathrm{adm}}}]=0.
\]
Thus locality survives admissibility because admissibility is implemented as a state-space
restriction, not as a nonlocal compression defining the algebra itself.
\end{proof}

\begin{figure}[t]
\centering
\scriptsize
\begin{tikzpicture}[
    >=Latex,
    selector/.style={draw=blue!60!black, fill=blue!6, rounded corners, text width=2.65cm, minimum height=1.0cm, align=center},
    dyn/.style={draw=orange!75!black, fill=orange!8, rounded corners, text width=2.65cm, minimum height=1.0cm, align=center},
    result/.style={draw=green!50!black, fill=green!8, rounded corners, text width=2.9cm, minimum height=1.0cm, align=center},
    note/.style={font=\footnotesize\itshape, text=black!60, align=center}
]
\node[selector] (padm) at (-4.8,1.15) {$\Padm$\\state-space selector};
\node[dyn] (zrel) at (-1.6,1.15) {$Z_{\mathrm{rel}},\ W_{\mathrm{rel}}$\\relative dynamics};
\node[dyn] (sch) at (1.6,1.15) {$\{S_n^T\}$\\Schwinger family};
\node[dyn] (os) at (4.8,1.15) {$(\mathcal H_{\mathrm{OS}},\mathfrak A_M)$\\OS/Wightman theory};

\node[selector] (hadm) at (-1.6,-1.2) {$\mathcal H_{\mathrm{adm}}=\operatorname{Ran}(\Padm)$\\physical sector};
\node[result] (aadm) at (2.6,-1.2) {$\mathfrak A_{\mathrm{adm}}(\mathcal O)$\\local net by restriction};

\draw[->, thick] (zrel) -- (sch);
\draw[->, thick] (sch) -- (os);
\draw[->, thick, dashed] (padm) -- (hadm);
\draw[->, thick] (os) -- (aadm);
\draw[->, thick] (hadm) -- (aadm);

\node[note] at (0,-2.35) {Locality is inherited by restriction to $\mathcal H_{\mathrm{adm}}$, not
by defining the net through a global compression $\Padm A \Padm$.};
\end{tikzpicture}
\caption{Selector versus dynamics on the admissible branch. The projector $\Padm$ selects the
physical state space, while Schwinger functions, OS reconstruction, and the later exact flow carry
the local Minkowski dynamics.}
\label{fig:selector-versus-dynamics}
\end{figure}

\begin{theorem}[\textnormal{[M]} Stable massive scattering closure on the admissible branch]
\label{thm:stable-massive-scattering}
Let $a$ be a stable massive species on $\mathcal U_{\mathrm{scat}}$ such that the renormalized
two-point function has an isolated simple pole
\[
\widetilde G_a^{(2)}(p)
=
\frac{Z_a}{p^2+m_a^2-i0}+R_a(p),
\qquad
Z_a>0,
\]
with $R_a(p)$ regular near $p^2=-m_a^2$ and $m_a>0$ below the first inelastic threshold. Then the
Haag--Ruelle asymptotic fields $\Phi_a^{\mathrm{in/out}}$ exist on $\mathcal H_{\mathrm{adm}}$, the
wave operators $\Omega_\pm$ exist on the stable massive scattering sector, and
\[
S_{\mathrm{massive}}:=\Omega_+^\ast\Omega_-
\]
is well defined there.
\end{theorem}

\begin{proof}
By \TFPTxref{thm:admissible-local-minkowski-net,thm:lorentzian-reconstruction}, the admissible branch
has a local relativistic net with positive Hamiltonian on $\mathcal H_{\mathrm{adm}}$. The isolated
simple pole produces the one-particle spectral projector of
\Cref{lem:one-particle-projectors-from-poles}, and exponential clustering on the gapped sector is
provided by \TFPTxref{thm:nonperturbative-admissible-qft}. The appendix theorem
\Cref{thm:haag-ruelle-admissible-massive} therefore applies and yields the asymptotic fields, wave
operators, and the stable massive scattering operator.
\end{proof}

\begin{theorem}[\textnormal{[M]} LSZ reduction on the stable massive sector]
\label{thm:lsz-stable-massive}
For stable massive external legs on the admissible branch, connected scattering amplitudes satisfy
\[
\mathcal A_{m\leftarrow n}^{\mathrm{conn}}
=
\Bigl(\prod_{r=1}^{m+n} Z_r^{-1/2}\Bigr)
\lim_{\mathrm{on\ shell}}
\Bigl(\prod_{r=1}^{m+n}(p_r^2+m_r^2)\Bigr)
\widetilde G_{m+n}^{\mathrm{conn}}(p_1,\ldots,p_{m+n}).
\]
Equivalently, the stable massive scattering matrix is reconstructed from the renormalized
admissible correlators of the theory.
\end{theorem}

\begin{proof}
By \Cref{thm:stable-massive-scattering}, the stable massive sector admits asymptotic fields and wave
operators. The appendix corollary \Cref{cor:lsz-admissible-massive} then expresses matrix elements
of $S_{\mathrm{massive}}$ in terms of amputated connected correlators with the displayed residue
factors.
\end{proof}

\begin{remark}[Resonant states]
\label{rem:resonant-states-admissible}
If a pole of the renormalized inverse propagator is complex,
\[
s_\ast=M_a^2-iM_a\Gamma_a,
\]
then no asymptotic one-particle vector is claimed. The physical content is the pole position and
the residue of the corresponding partial amplitudes, not an element of the asymptotic Hilbert
space.
\end{remark}

\begin{theorem}[\textnormal{[B]} Low-curvature dressed massless scattering interface]
\label{thm:dressed-massless-interface}
\label{rem:ir-dressed-massless-sectors}
On the low-curvature laboratory region $\mathcal U_{\mathrm{scat}}$, the photon and linearized
graviton sectors admit infrared-dressed M\o ller maps
\[
\Omega_\pm^{\mathrm{dress}}.
\]
The corresponding inclusive scattering operator
\[
S_{\mathrm{dress}}:=(\Omega_+^{\mathrm{dress}})^\ast \Omega_-^{\mathrm{dress}}
\]
is the correct asymptotic interface object on the massless branch. Bare undressed massless
amplitudes are not used as the unitarity object.
\end{theorem}

\begin{proof}
By \TFPTxref{thm:admissible-local-minkowski-net,thm:lorentzian-reconstruction}, the admissible branch
provides a local low-curvature asymptotic net carrying the photon and linearized graviton sectors.
In these long-range channels the correct asymptotic object is the infrared-dressed inclusive
scattering operator rather than the bare Fock-space $S$-matrix. Applying the standard long-range
dressing machinery to the admissible soft charges therefore yields the dressed M\o ller maps and
the interface operator $S_{\mathrm{dress}}$.
\end{proof}

\begin{definition}[Stable record algebra from the admissible gap]
\label{def:stable-record-algebra}
Let
\[
\tau_*:=\frac{1}{m_*}.
\]
For a tolerance $\varepsilon>0$, define
\[
\Arec(t)
:=
\left\{
O\in\mathfrak A_{\mathrm{adm}}
:\ \|[O(t),O(t')]\|\le\varepsilon
\text{ for } |t-t'|>\tau_*
\right\}.
\]
\end{definition}

\begin{proposition}[Stable records from gap and clustering]
Under \TFPTxref{thm:nonperturbative-admissible-qft}, the algebra $\Arec(t)$ defines the slow, robust sector singled out by the admissible mass gap. Record observables are therefore a dynamical consequence of gap plus clustering rather than an independent philosophical postulate.
\end{proposition}

\begin{proof}[Proof sketch]
The spectral gap suppresses long-time commutators by the same exponential scale that controls connected correlators. Therefore operators whose support is carried by the low-frequency admissible sector become approximately commuting on time separations larger than $\tau_*=1/m_*$, which is exactly the stability criterion built into \Cref{def:stable-record-algebra}.
\end{proof}

\begin{remark}[Record and empirical readout layers]
Stable-record algebras, appendix-level readout algebras, and Born-type evaluation are deferred to
Appendix~\TFPTxref{app:record}. They are not used in the main text to select the closed branch.
\end{remark}

\begin{remark}[Appendix horizon continuation]
Horizon applications are deferred to Appendix~\TFPTxref{app:horizons}. They are not part of the closed
main-text theorem chain.
\end{remark}

\begin{corollary}[Master closure corollary]
\label{cor:master-closure-corollary}
The roadmap items (1)--(7) of \Cref{prop:master-closure-roadmap} are now established by
the corresponding individual theorems of this section, with no item appearing as an assumption:
item (1) by \TFPTxref{thm:apsdom-upgrade}, item (3) by \TFPTxref{thm:padm-upgrade}, item (4) by
\TFPTxref{thm:topological-family-closure}, item (5) by the determinant-line consequence of the
master closure roadmap together with the family closure, item (6) by
\TFPTxref{thm:reflection-positivity-admissible}, and item (7) by the lower-bound and coercivity
clauses underlying \Cref{thm:admissible-vacuum-finite-volume}. Hence the roadmap reads as a
master closure corollary rather than as a hidden master assumption.
\end{corollary}


% ---- Internal reduction theorem and renormalized observable hierarchy: source lines 11328-12235 ----
